\section{Quantum Fermi Gas}
\label{sec:fermi}
Before we discuss neutron stars as Fermi gas, we must discuss Fermi gas. The formalism developed in this section serves two purposes: it provides the equation of state relating pressure to energy density in the neutron star interior, and it gives us the tools needed to assess whether the zero-temperature (cold gas) approximation is justified at the densities found within these objects.

\subsection{Fermi--Dirac Statistics and the Density of States}
 
Neutrons are spin-$\tfrac{1}{2}$ fermions, subject to the Pauli
exclusion principle: no two neutrons may occupy the same quantum state
simultaneously. In thermal and chemical equilibrium, the mean occupation
number of a single-particle state with energy $\epsilon$ is given by the
Fermi--Dirac distribution,
%
\begin{equation}
    f(\epsilon) = \frac{1}{\exp\!\left[(\epsilon - \mu)/k_{\mathrm{B}}T\right] + 1},
    \label{eq:fermi-dirac}
\end{equation}
%
where $\mu$ is the chemical potential, $k_{\mathrm{B}}$ is Boltzmann's
constant, and $T$ is the temperature. For a gas of $N$ neutrons in a
volume $V$, counting momentum states with periodic boundary conditions
and accounting for the two spin projections ($m_s = \pm\tfrac{1}{2}$),
the number of states in the momentum interval $[p,\, p + dp]$ is
%
\begin{equation}
    g(p)\,dp = \frac{8\pi V}{h^3}\,p^2\,dp,
    \label{eq:dos}
\end{equation}
%
where $h$ is Planck's constant. This density of states is the
fundamental object from which all thermodynamic quantities of the ideal
Fermi gas are derived.

\subsection{The Zero-Temperature Limit and the Fermi Energy}

 
At zero temperature the Fermi-Dirac distribution \eqref{eq:fermi-dirac}
reduces to a step function: all states with $\epsilon < \mu(T=0) \equiv
E_F$ are occupied, and all states above $E_F$ are empty. The quantity
$E_F$ is the Fermi energy, and the corresponding momentum
$p_F = \sqrt{2m_n E_F}$ (non-relativistically) is the Fermi
momentum. Applying the number constraint where $N = \sum_{k}f(\epsilon)$
at $T = 0$ using \eqref{eq:dos} gives $p_F$ in terms of $n$.

\begin{equation}
    n \equiv \frac{N}{V}
      = \frac{8\pi}{h^3}\int_0^{p_F} p^2\,dp
      = \frac{8\pi}{3h^3}\,p_F^3,
    \label{eq:number-density}
\end{equation}
\begin{equation}
    p_F = \left(\frac{3h^3 n}{8\pi}\right)^{1/3}.
    \label{eq:fermi-momentum}
\end{equation}

The corresponding Fermi temperature $T_F \equiv E_F / k_{\mathrm{B}}$
sets the scale below which quantum effects dominate. When $T \ll T_F$
the thermal smearing of the Fermi--Dirac distribution over an energy
range ${\sim}k_{\mathrm{B}}T$ near $E_F$ is negligible, and the
zero-temperature approximation is excellent. We will check if this approximation is appropriate when talking about Neutron Star formulation.
 
\subsection{Degeneracy Pressure}
 
The Pauli exclusion principle forces neutrons to occupy states of
progressively higher momentum as the density increases, contributing
kinetic energy and hence pressure even at zero temperature and in the
absence of inter-particle interactions. degeneracy pressure
is what supports a neutron star against gravitational collapse.
 
The total kinetic energy of the gas at zero temperature is

\begin{equation}
    E_{\mathrm{kin}}
      = \frac{8\pi V}{h^3}\int_0^{p_F} \mathcal{E}(p)\,p^2\,dp,
    \label{eq:total-energy}
\end{equation}
%
where $\mathcal{E}(p)$ is the single-particle energy. The pressure
follows from the thermodynamic relation $P = -(\partial E_{\mathrm{kin}}/\partial V)_N$,
yielding

\begin{equation}
    P = \frac{8\pi c}{3h^3}\int_0^{p_F}
        \frac{p^4\,dp}{\mathcal{E}(p)/c}.
    \label{eq:pressure-general}
\end{equation}

The specific form of $\mathcal{E}(p)$ determines which regime the gas
occupies.

\subsection{The Non-Relativistic Limit}
 
When $p_F \ll m_n c$, the rest mass energy dominates and the
single-particle energy is $\mathcal{E}(p) \approx m_n c^2 + p^2/2m_n$.
Evaluating \ref{eq:pressure-general} in this limit gives
%
\begin{equation}
    P_{non-rel}
      = \frac{1}{20}\left(\frac{3}{\pi}\right)^{2/3}
        \frac{h^2}{m_n}\,n^{5/3}
      \propto \rho^{5/3},
    \label{eq:eos-nr}
\end{equation}
%
where $\rho = m_n n$ is the mass density. This is a \emph{polytropic}
equation of state with index $\gamma = 5/3$, or equivalently polytropic
index $\mathfrak{n} = 3/2$ in the Lane--Emden formalism that will be
employed in Section~\ref{sec:newtonian}.
 
\subsection{The Ultra-Relativistic Limit}
 
When $p_F \gg m_n c$, the rest mass is negligible and the single-particle
energy approaches $\mathcal{E}(p) \approx pc$. Evaluating
\eqref{eq:pressure-general} in this limit gives
%
\begin{equation}
    P_{rel}
      = \frac{hc}{8}\left(\frac{3}{\pi}\right)^{1/3} n^{4/3}
      \propto \rho^{4/3}.
    \label{eq:eos-ur}
\end{equation}
%
The softening of the equation of state from $\gamma = 5/3$ to
$\gamma = 4/3$ has profound consequences for stellar stability. A
polytrope with $\gamma > 4/3$ can always find a new equilibrium if
slightly compressed, since the pressure rises faster than gravity.
When $\gamma = 4/3$ exactly, the star is marginally stable: no
equilibrium exists above a critical mass, regardless of how much the
central density increases. This is the physical origin of the maximum
mass for compact stars, and it is the argument that underlies the
Chandrasekhar limit for white dwarfs and, as shown in
Section~\ref{sec:newtonian}, a corresponding limit for neutron stars.
 
\subsection{The Full Relativistic Equation of State}
 
For neutron star interior conditions the Fermi momentum is
neither much less nor much greater than $m_n c$, so the full
relativistic dispersion relation $\mathcal{E}(p) = c\sqrt{p^2 + m_n^2 c^2}$
must be used. The energy density and pressure are then

\begin{equation}
    \varepsilon
      = \frac{8\pi c}{h^3}
        \int_0^{p_F} \sqrt{p^2 + m_n^2 c^2}\; p^2\,dp,
    \label{eq:energy-density-rel}
\end{equation}

\begin{equation}
    P = \frac{8\pi c}{3h^3}
        \int_0^{p_F}
        \frac{p^4\,dp}{\sqrt{p^2 + m_n^2 c^2}}.
    \label{eq:pressure-rel}
\end{equation}

These integrals interpolate smoothly between the non-relativistic result
\eqref{eq:eos-nr} and the ultra-relativistic result \eqref{eq:eos-ur} as
$p_F / m_n c$ increases from $0$ to $\infty$, and they constitute the
equation of state that Oppenheimer and Volkoff used in their original
calculation \cite{Oppy}. Closed-form expressions in terms of
standard functions exist \cite{Chandrasekhar1935} but the integrals are
straightforward to evaluate numerically.